\clearpage
\section{Design}

The solution uses a standard RMI split.
\begin{itemize}
    \item The registry only hosts the lookup service.
    \item The lobby is the matchmaking entry point and creates remote game objects.
    \item Each game instance owns one match, validates all moves, and keeps turn order and status consistent.
    \item Clients export callback objects and receive immutable \texttt{BoardState} snapshots.
\end{itemize}

The implementation keeps the code small by giving each class one
responsibility. \texttt{LobbyImpl} manages room lifecycle, \texttt{GameImpl}
manages game rules, and \texttt{GameControllerImpl} bridges remote callbacks
to the local UI or CLI. Finished matches are pruned after they reach a
terminal state.

The game model is intentionally snapshot-based. Clients receive immutable
\texttt{BoardState} objects, so remote updates never leak mutable server
internals. This makes the callback protocol easier to reason about and also
simplifies testing, because a received board can be compared directly without
worrying about later mutations.

\subsection{Trade-offs}

The implementation uses a clear server-centric design. It does not try
to introduce persistence or matchmaking policies beyond the assignment scope,
and it keeps the remote protocol small. The concurrency boundaries remain on
the server side, where game state, turn validation, and terminal conditions
are enforced.
